\color{unimain}{\subsection{Логика приема сигналов внешнего отключения}}
\color{black}

\begin{enumerate}

\item
Логика приема сигналов внешнего отключения от РЗ ВН выводится накладкой \textbf{«ЛО внеш РЗ ВН»}.

\item
Команда на срабатывание \textit{«ЛО внеш РЗ: Срабатывание»} формируется при возникновении любого из следующих событий:
\begin{itemize}
\item прием сигналов отключения от внешних устройств РЗ;
\item прием сигналов отключения от внешних УРОВ;
\item прием сигнала отключения от ТЗНП смежного трансформатора.
\end{itemize}

\item
Запрет АПВ от внешних устройств РЗ \textit{«ЛО внеш РЗ: Запрет АПВ»} формируются при возникновении любого из следующих событий:
\begin{itemize}
\item прием сигналов отключения от внешних устройств РЗ;
\item прием сигналов отключения от внешних УРОВ;
\item прием сигнала отключения от ТЗНП смежного трансформатора при включенной накладке \textbf{«ЗАПВ от ТЗНП смеж. Т»}.
\end{itemize}

\item
Алгоритм работы ЛО внеш РЗ ВН приведен на рисунке \ref{pic:LO_VN}. Параметры для настройки ЛО внеш РЗ ВН приведены в таблице \ref{app_set:lovneshvn}.

\medskip
\begin{figure}[H]
\centering
\includegraphics[width=1\textwidth,height=1\textheight,keepaspectratio]{\fbpath/06.01 ДЗТ/021.9901 ЛО внеш/img/LO_VN.pdf}
\captionof{figure}{Функциональная схема логики приема сигналов внешнего отключения стороны ВН}\label{pic:LO_VN}
\end{figure}

\item
Логика приема сигналов внешнего отключения от РЗ НН1(НН) и РЗ НН2(СН) выводится накладкой \textbf{«ЛО внеш РЗ НН1»} и \textbf{«ЛО внеш РЗ НН2»} соответственно.
 
\item
Команды на срабатывание \textit{«ЛО внеш РЗ НН1(НН): Срабатывание»} и \textit{«ЛО внеш РЗ НН2(СН): Срабатывание»} формируются при прием сигналов отключения от внешних устройств РЗ (ДЗО, ЗДЗ, УРОВ).

\item
Алгоритм работы ЛО внеш РЗ НН приведен на рисунке \ref{pic:LO_SN}. Параметры для настройки ЛО внеш РЗ НН приведены в таблице \ref{app_set:lovneshnn}.

\medskip
\begin{figure}[H]
\centering
\includegraphics[width=1\textwidth,height=1\textheight,keepaspectratio]{\fbpath/06.01 ДЗТ/021.9901 ЛО внеш/img/LO_NN.pdf}
\captionof{figure}{Функциональная схема логики приема сигналов внешнего отключения сторон СН и НН}\label{pic:LO_SN}
\end{figure}

\item
Логика приема сигналов внешнего отключения от АПТ и АОТ выводится накладкой \textbf{«ЛО внеш~АПТ»}.

\item
Команда на срабатывание \textit{«ЛО внеш АПТ/ ЛО внеш АПТ: Срабатывание»} формируется при приеме сигналов отключения от внешних устройств АПТ или АОТ.

\item
Алгоритм работы ЛО внеш АПТ приведен на рисунке \ref{pic:LO_APT}. Параметры для настройки ЛО внеш АПТ приведены в таблице \ref{app_set:lovneshapt}.

\medskip
\begin{figure}[H]
\centering
\includegraphics[width=1\textwidth,height=1\textheight,keepaspectratio]{\fbpath/06.01 ДЗТ/021.9901 ЛО внеш/img/LO_APT.pdf}
\captionof{figure}{Функциональная схема логики приема сигналов внешнего отключения от АПТ и АОТ}\label{pic:LO_APT}
\end{figure}


\end{enumerate}